\subsection{Module JajaCode : Compilation et Interprétation}

La réalisation du module JajaCode s'articule autour de deux responsabilités distinctes : la génération de code (depuis le MiniJaja) et l'exécution virtuelle. Voici les choix architecturaux qui ont guidé cette implémentation.

\subsubsection{Stratégie de Compilation : Externalisation de la logique}

Pour la phase de compilation (traduction du MiniJaja vers le JajaCode), nous avons écarté l'idée intuitive d'intégrer une méthode compile() directement au sein des classes de nœuds de l'AST MiniJaja.\\

L'utilisation stricte du Pattern Visiteur pour externaliser la logique de compilation.  Ce choix répond au principe de responsabilité unique (\textit{Single Responsibility Principle}). Il permet de conserver des modèles de données (les nœuds) purs et légers, tout en isolant la complexité de la génération d'instructions (gestion des adresses, étiquettes) dans une classe dédiée. Cela facilite également l'évolution du compilateur sans risquer de briser la structure du langage source.

\subsubsection{Stratégie d'Interprétation : Pattern Command et Refactoring}

L'implémentation initiale de l'interpréteur reposait sur une structure monolithique (une classe unique de plus de 1200 lignes avec un immense switch/case), ce qui rendait la maintenance et le débogage extrêmement ardus.\\

Nous avons opéré une refonte complète de l'interpréteur en appliquant le \textit{Pattern Command}. Chaque instruction JajaCode (ex: Add, Push, Invoke) est désormais encapsulée dans sa propre classe (nommée "Axiome" dans notre architecture) qui implémente une interface commune d'exécution.\begin{itemize} \item \textbf{Maintenabilité :} Cela fragmente la complexité. Chaque classe est responsable d'une seule instruction, rendant le code lisible. \item \textbf{Extensibilité :} L'ajout d'une nouvelle instruction ne demande pas de modifier le moteur central, mais simplement de créer une nouvelle classe. \end{itemize}

\subsubsection{Encapsulation du contexte d'exécution}

Pour orchestrer l'exécution de ces commandes, nous avons introduit une classe MachineContext.\\

La centralisation de l'état machine.Plutôt que de passer des références éparses à la Pile, au Tas à chaque instruction, l'objet MachineContext encapsule l'intégralité de la mémoire et de l'adressage. Les instructions interagissent uniquement avec ce contexte, assurant une forte cohésion des données.

\subsubsection{Gestion des erreurs et Testabilité}

Enfin, chaque "Axiome" a été conçu pour lever des exceptions spécifiques en cas d'échec (ex: StackUnderflowException, TypeException). Ce choix a été piloté par la stratégie de test. En typant fortement les exceptions, nous avons pu écrire des tests unitaires précis pour chaque instruction, validant non seulement le succès des opérations mais aussi leur robustesse face aux cas limites.

\subsection{Interface}
\subsubsection{Choix de réalisation}
Pour l’interface, nous avons décidé de créer un IDE qui ressemble à ceux de JetBrains, avec lesquels nous sommes familiers.

Nous avons choisi d’utiliser le composant CodeArea de RichTextFX, car il n’était pas disponible dans la version de JavaFX utilisée. Nous avons également choisi de ne pas monter en version de JavaFX afin de rester cohérents avec la version de Java utilisée et d’avoir la même version partout, car les versions plus récentes de JavaFX nécessitaient une version plus récente de Java.

Nous avons aussi choisi de ne pas utiliser les fichiers FXML pour créer notre interface, afin d’avoir une plus grande liberté et davantage de contrôle sur la structure de l’interface. Cela nous a permis de construire l’interface petit à petit et d’ajouter de nouveaux éléments sans avoir à tout refaire.

\subsubsection{Gestion des erreurs}
Côté interface graphique, nous avons l’affichage des erreurs de compilation et d’exécution dans la console. Nous disposons également de la coloration syntaxique, qui affiche en rouge les erreurs de syntaxe.